\section{Introduction}
\label{sec:intro}

A coding agent is a model inside a harness.
The harness owns the agent loop: it builds model requests, executes tool calls, returns results to the model, and manages context.
Claude Code, Codex CLI, OpenCode, Pi and fx are harnesses in this sense.
Benchmarks such as SWE-bench~\citep{jimenez2024swebench} and Terminal-Bench~\citep{merrill2026terminalbench} measure whether an agent completes a task when its tools work.
In practice tools do not always work.
A command times out, a file read returns nothing, a service answers with data that is out of date, or a call that took effect reports an error.

What the agent does next decides the outcome.
It can retry, take another route to the same information, stop, or continue as if nothing happened.
It can also report success for work that failed.
These behaviors depend on the model and on the harness: how the harness renders an error, whether it retries on its own, what it shows the model of a failed command, and how it lets the agent recover.

CMCA injects controlled, seeded faults into the tool calls of coding-agent harnesses and compares each configuration against its own clean baseline on the same tasks.
It addresses three questions, each mapped to metrics in Section~\ref{sec:metrics}:
\begin{itemize}
\item \textbf{RQ1, robustness.} How much does task success drop under faults, and how does each harness configuration compare when the main model is held fixed?
\item \textbf{RQ2, honesty.} How often does the agent report success for work that failed, or carry stale data into its output?
\item \textbf{RQ3, recovery.} What does the agent do after a fault, and what does recovery cost in tokens and steps?
\end{itemize}

This paper makes four contributions.
\begin{enumerate}
\item A fault model for coding-agent tools (Section~\ref{sec:faults}): fault types delivered before or after execution, a deterministic candidate draw keyed on the call, and persistence modes that keep every task solvable.
\item A two-layer injection design (Section~\ref{sec:design}): an MCP proxy that gives every harness the same bytes, and native injectors that reach each harness's built-in tools, with comparisons restricted to fault types a conformance suite verifies per harness.
\item Metrics beyond success (Section~\ref{sec:metrics}): overclaiming, propagation of stale data, duplicate side effects, recovery actions per fault event, and cost, with hierarchical bootstrap intervals and paired tests.
\item An open implementation on Harbor~\citep{harbor} (Section~\ref{sec:impl}) and a first measurement of five configurations per track with one model (Section~\ref{sec:results}), with every trial's trajectory and fault log.
\end{enumerate}
The results are descriptive: they expose behavior that success alone hides, such as frequent overclaiming, but the task suite is too small to establish a statistically supported ranking of harnesses.
